\begin{abstract}
Mixture-of-Experts (MoE) and dynamic adapter-merging systems have emerged as powerful paradigms for scalable model serving on edge devices. However, when deployed under sequential, heterogeneous query streams, serving controllers face a fundamental \emph{accuracy-stability dilemma}, characterized by the \emph{routing jitter paradox}. Stateless routers adapt quickly but suffer from high-frequency spatial (layer-to-layer) oscillations of merging weights, causing catastrophic representation drift. Conversely, stateful routers filter spatial jitter using temporal biochemical kinetics or differential moving averages, but introduce severe inertial lag (hysteresis) during rapid task switches, leading to downstream accuracy collapse. 

In this work, we propose \textbf{Markovian Path-Integral Ensembling (QPathMerge)}, a novel, training-free serving controller that completely resolves this dilemma. Inspired by statistical physics and path-integral formulations, we model the sequence of network layers as a discrete 1D lattice and represent the routing trajectory as a discrete Euclidean path integral over network depth. By mapping this formulation to a Markov Random Field, we execute the \textbf{Forward-Backward sum-product algorithm (Belief Propagation)} to calculate the exact, globally optimized marginal probability of expert selection at each layer in $O(L K^2)$ time. Because QPathMerge propagates messages symmetrically across network depth for each input sample independently, it acts as a mathematically optimal spatial low-pass filter (achieving near-oracle trajectory smoothness) while maintaining absolute statelessness across sequence samples (eliminating temporal serving lag). Evaluation inside a high-fidelity 14-layer Coordinate Sandbox shows that QPathMerge slashes layer-wise routing jitter by over $3\times$ compared to state-of-the-art stateful kinetics models, while completely avoiding representation hysteresis to achieve a leading \textbf{98.50\%} heterogeneous serving accuracy.
\end{abstract}
